%
%	Configure LaTeX to produce a PDF presentation using the beamer class
%


% Create the slide deck itself
% Default is 16:9 widescreen ratio
\documentclass[ignorenonframetext,12pt,aspectratio=169]{beamer}


% Create a handout with the slides on a page
% \documentclass[handout]{beamer}
% \usepackage{pgfpages}
% \usepackage{handoutWithNotes}
% % \pgfpagesuselayout{4 on 1}[letterpaper,landscape,border shrink=5mm]
% \pgfpagesuselayout{3 on 1 with notes}


% Allow things like '<a>' in verbatim environments (fixes font encoding issue)
\usepackage[T1]{fontenc}


% Required packages to suppport MMD features
\usepackage[sort&compress]{natbib}  % Support bibliographies
\usepackage[acronym]{glossaries}    % Support glossary entries
\usepackage{booktabs}               % Better tables
\usepackage{tabulary}               % Support longer table cells

% Support CriticMarkup features
\usepackage[normalem]{ulem}         % Support strikethrough
\usepackage{soul}
\usepackage{xargs}
\usepackage{todonotes}
\presetkeys{todonotes}{inline}{}    % We need inline notes in Beamer
\newcommandx{\cmnote}[2][1=]{\linespread{1.0}\todo[linecolor=red,backgroundcolor=red!25,bordercolor=red,#1]{#2}}

% Use \ul instead of \underline since we are using soul
\let\underline\ul


% For each part, show a title page and an outline of that part
\AtBeginPart
{
   \begin{frame}
      \partpage
   \end{frame}

   \begin{frame}
       \frametitle{Outline}
       \tableofcontents
   \end{frame}
}

% For each section, show an outline highlighting where we are
\AtBeginSection[]
{
   \begin{frame}
       \frametitle{Outline}
       \tableofcontents[currentsection,currentsubsection]
   \end{frame}
}
